% ch07 —— §7 穿墙公式的一个证明 + 致谢（原文 wcf-paper.tex 3153–3239 行，PDF 第 50–51 页）
\section{穿墙公式的一个证明}
\label{sec:proof}

本文中，我们在 $\CM$ 上构造了一个 \hk 度量 $g$，并论证它与 $\IR^3\times S^1$ 上该规范理论的模空间所带的物理度量相符。Kontsevich--Soibelman 穿墙公式是作为一条自洽性条件出现的：没有它，我们的构造就无法给出光滑的度量。我们把这看作穿墙公式确实成立的有力旁证。

然而，这些构造还称不上是对穿墙公式的一个证明。要给出证明，我们必须直接从规范理论的物理出发，而不对度量应当采取何种形式作任何假设。在这种进路下，我们无法动用黎曼--希尔伯特构造这一工具（至少一开始不能）。我们改用第 5.5 节所描述的另一种视角。设 $\CM$ 是该规范理论在 $\IR^3\times S^1$ 上的模空间，并考虑映射 $\CX(\zeta)\colon \CM\to T$（$\zeta\neq0,\infty$）。我们的目标是为此类 $\CX(\zeta)$ 构造一组可积方程，其形式如 (5.30)–(5.35)，使得 $\IC\IP^1$ 上的联络 (5.35) 具有斯托克斯射线 $\ell$，其上承载斯托克斯因子
\begin{equation} \label{eq:sf-redux}
S_\ell = \prod_{\gamma \in (\clat_u)_\ell} \CK_\gamma^{\degen(\gamma;u)}.
\end{equation}
一旦构造出这样的微分方程，WCF 便是说：联络 (5.35) 经历等单值形变。% TODO 合稿改 \ref{sec:differential-problem}、\eqref{eq:conn-holomorphy-1}–\eqref{eq:zeta-equation}

下面我们说明如何直接从规范理论导出这些微分方程。正如我们在附录 B 中所说明的，(5.30)、(5.31) 有一个简单的几何意义：它们正是柯西--黎曼方程，表达 $\CX(\zeta)$ 关于复结构 $J^\pz$ 的全纯性。特别地，这些方程可以纯粹用约化理论的 $\CN=4$ 超对称来理解。其次注意到，(5.32)、(5.33) 与 (5.30)、(5.31) 的形式完全相同。事实上，只要把 $\Lambda$ 看作某个「背景」矢量多重态的标量分量，它们就会变得完全一样。这是证明非重整化定理的一个标准技巧，例如见 \cite{Seiberg:1993vc}；把它用在这里，应当就能得到所需的 (5.32)、(5.33)。若理论含有质量参数，我们还可以通过弱规范化味对称性，类似地证明 (6.11)、(6.12)。% TODO 合稿改 \ref{app:holomorphy}、\eqref{eq:conn-holomorphy-1}、\eqref{eq:conn-holomorphy-2}、\eqref{eq:lambda-equation}、\eqref{eq:lambda-bar-equation}、\eqref{eq:mass-equation}、\eqref{eq:mass-bar-equation}

最后，我们需要建立关键方程 (5.34)、(5.35)，它们给出 $\CX$ 对 $R$ 与 $\zeta$ 的依赖。这两个方程可从 $d=4$ 理论的反常 $U(1)_R$ 对称性与标度不变性推出，而这些对称性由 (5.42)、(5.43) 或 (6.13)、(6.14) 表达；把上述对称性与上文已经建立的方程结合起来，即可得到它们。诸函数 $\CX_\gamma$ 有一种物理解释，我们希望另文描述。可以把它们看作某个三维拓扑场论的手征环的元素，或看作四维理论中某些线算子的期望值。按这种眼光，(5.34)、(5.35) 就是反常 Ward 恒等式。% TODO 合稿改 \eqref{eq:R-equation}、\eqref{eq:zeta-equation}、\eqref{eq:rescaling-1}、\eqref{eq:rescaling-2}、\eqref{eq:rescaling-m1}、\eqref{eq:rescaling-m2}

为完成证明，我们还须说明：$\IC\IP^1$ 上该联络的斯托克斯因子确实由 \eqref{eq:sf-redux} 给出。为此我们利用这样一件事实：由 (5.34)，斯托克斯因子在 $R$ 变化下保持不变。于是我们可以取非常大的 $R$；在那里（远离诸墙处），度量的修正应当可以很好地近似为若干单瞬子修正的线性叠加，而这些单瞬子修正是我们在阿贝尔理论中已经知道的。从 $\IC\IP^1$ 上的联络转到相应的黎曼--希尔伯特问题，并重复我们在第 5 节用过的论证，便可说明这些修正直接对应于斯托克斯因子。至此，穿墙公式的证明完成——至少是在物理的严格性水平上。% TODO 合稿改 \eqref{eq:R-equation}、\ref{sec:constructing-moduli}

\section*{致谢}

我们要特别感谢 F. Denef、N. Nekrasov 与 X. Yin 在本项目早期阶段的合作，并感谢 B. Pioline 关于半平坦度量的重要讨论。我们也感谢 Wu-yen Chuang、E. Diaconescu、D. Jafferis、D. Joyce、M. Kontsevich、S. Lukyanov、N. Seiberg、Y. Soibelman、V. Toledano Laredo、C. Vafa、E. Witten 与 E. Zaslow 的有益讨论。GM 的工作受美国能源部（DOE）资助，资助号 DE-FG02-96ER40949。GM 还感谢阿斯彭物理中心（Aspen Center for Physics）与加州大学圣巴巴拉分校（UCSB）的卡弗里理论物理研究所（KITP）在本工作完成期间的热情接待；KITP 的部分经费来自美国国家科学基金会（NSF）的资助 PHY05-51164。AN 的工作部分由高等研究院（Institute for Advanced Study）的 Martin A. and Helen Chooljian 会员资格资助，部分由 NSF 在资助号 PHY-0503584 下资助。DG 的工作部分由 DOE 资助号 DE-FG02-90ER40542 资助。

\begin{yizhu}
本节中的 $T$ 指 $\CM$ 的环面纤维（$\CM$ 在拓扑上是 4d 模空间 $\vm$ 上的 \SW 环面丛），映射 $\CX(\zeta)$ 取值于其中；$\CK_\gamma$ 是第 2.2 节引入的 Kontsevich--Soibelman 辛同胚，$(\clat_u)_\ell$ 则是中心荷相位落在射线 $\ell$ 上的那些电荷所构成的子集。% TODO 合稿改 \ref{sec:ReviewKS}
\end{yizhu}
